% GENERATED FILE. Edit canonical lessons, then run npm run generate:books.
% canonical-source-hash: fnv1a64:7a5542ec6da59f5d
% canonical-lessons: IT-C270-piu-di, IT-C270-da-tre-anni, IT-C270-ripasso-esperienza

\chapter{More Than, For Three Years --- Comparing and How Long}
\label{ch:it-comparing-and-how-long}
\hlchaptermodality{270}

\begin{chapteropening}
\textbf{By the end of this chapter:} \emph{I can compare two things with più … di, name the most with il più, and say how long with da and per.}

Complete IT-C270-ripasso-esperienza: describe an experience with a feeling, a hope, a comparison and a length.
\end{chapteropening}

\section[più … di]{più … di — more than}

\label{lesson:IT-C270-piu-di}



\begin{quote}
Say \textquotedblleft{}see you later\textquotedblright{}. (\emph{A più tardi.}) Literally, \textquotedblleft{}to more late\textquotedblright{}.

That \textbf{più} is the whole of this lesson's new grammar.
\end{quote}

\subsection*{You'll want to know: più … di}
To compare two things, put \textbf{più}, \textquotedblleft{}more\textquotedblright{}, in front of the adjective and \textbf{di}, \textquotedblleft{}than\textquotedblright{}, in front of the second thing:

\begin{quote}
Il treno è \textbf{più} veloce \textbf{dell'}autobus. — The train is faster than the bus. L'aereo è \textbf{più} caro \textbf{del} treno. — The plane is dearer than the train. Sei \textbf{più} alto \textbf{di} me. — You are taller than me.
\end{quote}

English sometimes builds the comparison into the adjective, \emph{faster}, \emph{taller}; Italian builds it the same way almost every time, with \emph{più}. The \textbf{di} welds to the article, as \emph{a} does: \textbf{del}, \textbf{della}, \textbf{dell'}. Before a name or a stressed pronoun it stays plain: \emph{di me}, \emph{di te}.

Its opposite is \textbf{meno}, \textquotedblleft{}less\textquotedblright{}: \emph{L'autobus è \textbf{meno} caro \textbf{del} treno} — the bus is less expensive than the train.

\begin{grammarlens}[title={Grammar Lens: il più}]
For the most of all, keep \emph{più} and give the noun its article:

\begin{quote}
\emph{il treno \textbf{più} veloce} — the fastest train \emph{la città \textbf{più} bella} — the most beautiful city
\end{quote}

The adjective still comes after the noun, as Italian adjectives mostly do, and the article does the work of English \emph{-est}. Who it is the most among comes with \emph{di}: \emph{il treno più veloce \textbf{del} mondo}, the fastest train in the world.
\end{grammarlens}

\begin{cousinweb}
\textbf{più} is inherited from Latin \emph{plūs}, \textquotedblleft{}more\textquotedblright{} (EtymDB 2.1). Latin \emph{pl} became Italian \emph{pi}, the same change that made \emph{plicāre} into \textbf{piegare}. English borrowed \emph{plūs} whole as \textbf{plus}, and built on it \textbf{plural} and \textbf{surplus}.
\end{cousinweb}

\subsection*{Your turn}
Say these aloud:

\begin{itemize}
\raggedright
  \item \emph{Il treno è più veloce dell'autobus.}
  \item the other way round — \emph{L'autobus è meno veloce del treno.}
  \item \emph{Sei più alto di me.}
  \item \emph{il treno più veloce}, then \emph{la città più bella}
\end{itemize}

\begin{tcolorbox}[breakable,colback=teal!4,colframe=teal!35!black,title={Before you move on}]
Say \textquotedblleft{}the plane is dearer than the train\textquotedblright{}. (\emph{\textbf{L'aereo è più caro del treno.}}) How do you say \textquotedblleft{}the fastest train\textquotedblright{}? (\emph{\textbf{Il treno più veloce.}}) Which Latin word is \emph{più}, and which English word kept it whole? (\emph{\textbf{Plūs}} — \textbf{plus}.)
\end{tcolorbox}
\section[da tre anni]{da tre anni — how long, and still going}

\label{lesson:IT-C270-da-tre-anni}



\begin{quote}
Say \textquotedblleft{}the train is faster than the bus\textquotedblright{}. (\emph{Il treno è più veloce dell'autobus.}) Then say \textquotedblleft{}a year\textquotedblright{}. (\emph{Un anno.})

An experience has a length. Italian measures it two ways, and the choice depends on whether it is over.
\end{quote}

\subsection*{You'll want to know: da tre anni}
\begin{quote}
\textbf{Studio} l'italiano \textbf{da tre anni}. — I have been studying Italian for three years. \textbf{Abito} qui \textbf{da due mesi}. — I have lived here for two months.
\end{quote}

Listen to the verbs. English says \emph{have been studying}, a form that reaches back into the past. Italian says \textbf{studio}, the plain \textbf{present}, because the studying is still going on today. \textbf{da} with a length of time means \textquotedblleft{}from that far back, up to now\textquotedblright{}.

That is the one to practise until it is automatic, because an English speaker's instinct sends them to the past, and the past here says something different.

\begin{grammarlens}[title={Grammar Lens: per, for a stretch that is over}]
When the stretch has \textbf{ended}, the verb goes into the passato prossimo and the length takes \textbf{per}:

\begin{quote}
\textbf{Ho abitato} a Roma \textbf{per due anni}. — I lived in Rome for two years. \emph{(and then I left)}
\end{quote}

So one English word, \emph{for}, splits in two. \textbf{da} + present: it is still going on. \textbf{per} + passato prossimo: it is over. And \textbf{durante}, which you have, is \textquotedblleft{}during\textquotedblright{}: it names the stretch something happened inside, \emph{durante il viaggio}, \textquotedblleft{}during the journey\textquotedblright{}, not its length.
\end{grammarlens}

\begin{cousinweb}
\textbf{da} is two Latin prepositions run together: \emph{dē}, \textquotedblleft{}from, down from\textquotedblright{}, and \emph{ab}, \textquotedblleft{}from, away from\textquotedblright{} (EtymDB 2.1 lists it as a compound of the two). So \emph{da tre anni} is, quite literally, \textquotedblleft{}from three years\textquotedblright{}: counting from a point three years back, up to the present moment.
\end{cousinweb}

\subsection*{Your turn}
Say these aloud:

\begin{itemize}
\raggedright
  \item \emph{Studio l'italiano da tre anni.}
  \item \emph{Abito qui da due mesi.}
  \item the stretch that is over — \emph{Ho abitato a Roma per due anni.}
  \item your own — something you have been doing \emph{da} … and something you did \emph{per} …
\end{itemize}

\begin{tcolorbox}[breakable,colback=teal!4,colframe=teal!35!black,title={Before you move on}]
Which tense goes with \emph{da tre anni}, and why? (The \textbf{present} — it is still going on.) Which word takes over when the stretch is finished? (\emph{\textbf{Per}}, with the passato prossimo.) And what is \emph{da}, taken apart? (\emph{\textbf{Dē}} and \emph{\textbf{ab}}, \textquotedblleft{}from\textquotedblright{} twice.)
\end{tcolorbox}
\section[(review)]{Review — describe an experience}

\label{lesson:IT-C270-ripasso-esperienza}



\begin{quote}
Say \textquotedblleft{}I have been studying Italian for three years\textquotedblright{}. (\emph{Studio l'italiano da tre anni.})
\end{quote}

\begin{grammarlens}[title={What you've built}]
Two chapters gave you the four things that turn a fact into an \textbf{esperienza}: how it felt, what you hope, how it compared, and how long it went on.

\begin{quote}
\emph{Studio l'italiano da tre anni. Prima era difficile ed ero nervoso, ma adesso sono contento. Secondo me l'italiano è più facile della matematica. Spero di parlare bene.}
\end{quote}

Count the moves: a length still running, \emph{da tre anni}; a feeling in the past, \emph{ero nervoso}, against one now, \emph{sono contento}; a comparison, \emph{più facile della matematica}, marked as an opinion; a hope, \emph{spero di}.
\end{grammarlens}

\subsection*{Your turn}
Say these aloud:

\begin{itemize}
\raggedright
  \item \emph{Studio l'italiano da tre anni.}
  \item \emph{Prima era difficile ed ero nervoso, ma adesso sono contento.}
  \item \emph{Secondo me l'italiano è più facile della matematica.}
  \item a finished stretch — \emph{Ho abitato al mare per un anno, ed ero felice.}
  \item your own experience, with all four moves, ending on \emph{Il mio sogno è …}
\end{itemize}

\begin{tcolorbox}[breakable,colback=teal!4,colframe=teal!35!black,title={Before you move on}]
Name the four moves of an experience. (How it \textbf{felt}, what you \textbf{hope}, how it \textbf{compared}, how \textbf{long}.) Which tense goes with \emph{da}? (The \textbf{present}.) Say \textquotedblleft{}the most beautiful city\textquotedblright{}. (\emph{\textbf{La città più bella.}})
\end{tcolorbox}
